\chapter{Periods and Lattices}

The first, and arguably most important, mathematical foundation for this paper would be understanding the abstract concept of a lattice, and how the shortest vector in a lattice inherently quantifies the notion of period.

\textbf{Definition 5.1 (Lattice)} A lattice is a discrete additive subgroup of $\mathbb{R^n}$, i.e. it is a subset $\Lambda \subseteq \mathbb{R}$ satisfying the following properties:

1. $\Lambda$ is closed under addition and subtraction.\\
2. $\exists \:\epsilon > 0$ such that for any two distinct lattice points $x\neq y \in  \Lambda$, their distance is at least $\lVert x - y \rVert \geq \epsilon$.

In practice, this implies that a n-dimensional lattice $\Lambda^n$ is the set $L$ of all integer linear combinations of n linearly independent vectors. In other words:

\begin{equation}
    \Lambda = \set{z_1b_1+z_2b_2+\dots +z_nb_n | z_i \in \mathbb{Z},b_i \in \mathbb{R^n}}
\end{equation}
The vectors $b_i$ form a basis of the lattice. The most common example of a lattice is $\mathbb{Z}^n$, the set of all n dimensional vectors with integer entries. Another well-known lattice, which is the one that will concern us in this paper, is the set of all vectors satisfying a linear modular multiplicative congruence (if that isn't obvious from the description, those vectors need to have integer entries).

Now, given a lattice $\Lambda$, the Shortest Vector Problem (SVP) asks to find the non-zero vector with minimal Euclidean length. That is, we are looking for $v_0$, such that:
\begin{equation}
    \lVert v_0\rVert \leq \lVert v\rVert \:\forall \:v \in \Lambda
\end{equation}

This problem is, in general, NP-Hard\cite{Kreuzer2025}.

Now, the Period Finding problem (PF) can be stated as follows:\\
Given a function $f: \mathbb{Z}\to S$ that is \textbf{periodic}, meaning there exists some $r > 0$ s.t. $f(x+r) = f(x)\: \forall \:x$, find the smallest possible $r$ that satisfies this equation.

We will now show that, in 1 dimension, SVP reduces to period finding.

\textbf{Lemma 5.1 ($SVP \le_P PF$)} Define the set of periods $P$ as follows:
\begin{equation}
    P = \set{k\in \mathbb{Z}\: | \:f(x+k) = f(x)\: \forall \:x}
\end{equation}

\textit{Proof:} For $a,b \in P$,we have $f(x+a) = f(x)$, $f(x+b) = f(x)$. Therefore $f(x+ (a + b)) = f((x+a) + b) = f(x+a) = f(x)$, so $P$ is closed under addition.\\
Furthermore,$f(x-a) = f((x-a) + a) = f(x)$, meaning that if $a\in P$,then $-a \in P$, so every element has an inverse. Hence,since $P$ is closed under addition, it is also closed under subtraction. That proves the first lattice property for $P$.\\
Since $r$ is by definition the smallest element of $P$, setting $\epsilon = r$ satisfies property 2, and therefore $P$ is a lattice.

Moreover, since $P$ is a subset of $\mathbb{Z}$ that is closed under addition and every element has an inverse, it follows that $P$ is a subgroup of $\mathbb{Z}$. Every subgroup of $\mathbb{Z}$ is of the form $m\mathbb{Z} = \set{km : k \in \mathbb{Z}}$,where $m$ is the smallest element of the subgroup, so in this case $m = r$. Therefore, the elements of the lattice are $...,-2r,-r,0,r,2r,...$, and so the shortest nonzero elements are $r$ and $-r$, meaning the magnitude of the shortest vector in this case is exactly $r$. This concludes our proof.

This can easily be generalized to n dimensions. Suppose we have a function $f:\mathbb{Z}^n \to R$, which is periodic with respect to some arbitrary n-dimensional vector $v$. As stated above, a lattice can be constructed that consists of all other possible vectors with respect to which f is periodic. However, now we are no longer looking to find the shortest vector according to Euclidean distance, since that would simply be a vector with 0 everywhere besides the dimension with the shortest period. Instead,we want to find the shortest period in \textbf{every} dimension, so the shortest \textbf{basis} with respect to Euclidean distance. Therefore, the problem of finding periods in n dimensions broadly corresponds to finding the shortest basis of a given n-dimensional lattice.

To explaining the meaning of the word "broadly" there, we first need to define what a coset is\\

\textbf{Definition 5.2 (Coset)} Given some group $G$ and some some subgroup $H \subseteq G$, a coset with representative $g$ for some given $g \in G$ is simply the set $H_g$ defined as $H_g = \set{g + h\: \rvert h \in H}$.
 
We can think of $H_g$ as being essentially H shifted by g.

In reality, both SVP and factoring are special cases of a more general problem called HSP (Hidden Subgroup Problem) for specific types of groups \footnote{Read more about HSP in the \href{https://en.wikipedia.org/wiki/Hidden_subgroup_problem}{wikipedia page}}. Regev himself proved in \cite{Regev2004} that a quantum computer can solve SVP to some approximation factor (that is, find a vector that whose magnitude is larger than the shortest to a constant factor) under the assumption that there exists an algorithm that solves the hidden subgroup problem on the dihedral group by coset sampling . Factoring corresponds to solving HSP on finite abelian groups. In fact, the classical postprocessing in Regev's algorithm (most of what we will describe in Chapters 9 and 10) is essentially already outlined in the paper above - Regev's new idea is to create a relationship that intertwines factoring and SVP, solve SVP with coset sampling and use his solution to solve factoring as well.

Next, it is important to prove the following statement:

\textbf{Lemma 5.2} The set of vectors $V$ with entries $v_i \in \set{0,1\dots N-1}$ such that for a fixed set of bases $b_1,b_2\dots b_d$ it holds that
\begin{equation}
    \prod_{i=1}^d b_i^{v_i} = 1 \pmod N
\end{equation}

forms a subgroup of $\mathbb{Z}^d$ under vector addition, and therefore a lattice.\footnote{Eq 5.4 is a linear congruence in disguise. To see that, consider what happens when taking the logarithm of both sides.} 

\textit{Proof:} Clearly, $V \subseteq \mathbb{Z}^d$. Consider the requirements for $V$ to be a group:\\

Closure: $\prod_{i=1}^d b_i^{v_i+w_i} \pmod N = \prod_{i=1}^d b_i^{v_i} \cdot \prod_{i=1}^d b_i^{w_i} = 1 \cdot 1 = 1\pmod N$, hence if $b,w \in V$ then $b+w \in V$

Inverse element: $\prod_{i=1}^d b_i^{-v_i} = 1\pmod N$ since it is the modular inverse of the original product, which is congruent to 1 modulo N.

Identity element: setting $\mathbf{v = 0}$ satisfies the relationship.

Associativity is inherited by the parent group.

Therefore, the set of all vectors whose coordinates satisfy \textbf{Eq 5.4} forms a finite lattice $\Lambda$, and our proof is complete.

 Consider that a finite n-dimensional lattice modulo N whose elements are vectors that satisfy the congruence in \textbf{Eq 5.4} is a subgroup of the group of all vectors with integer coordinates mod N (denoted $Z_N^n$), so we can define a coset of that lattice with respect to that group accordingly. Crucially, replacing 1 with $u$ for some arbitrary $u \in \mathbb{Z}_N$ defines a coset whose representative is some exponent vector $\mathbb{t}$ that represents the exact shift that needs to be applied to the $v_i$ to get $u$ instead of 1.This will come into importance later.
 
Lastly, in our next chapters we will be discussing a lot the notion of a dual lattice. Let us define that here:\\

\textbf{Definition 5.3 (Dual lattice)} The dual lattice $\Lambda^*$ of a given lattice $\Lambda$, as follows:
\begin{equation}
    \Lambda^* = \set{\mathbf{y}\in \mathbb{R}^n \: | \: \langle \mathbf{y,x} \rangle \in \mathbb{Z} \quad \forall \: x \in \Lambda}
\end{equation}
where $\langle \mathbf{y,x} \rangle$ denotes the inner product operation between the $\mathbf{y}$ and $\mathbf{x}$ vectors.

